\documentclass[12pt,letterpaper]{article}
\usepackage{graphicx,textcomp}
\usepackage{natbib}
\usepackage{setspace}
\usepackage{fullpage}
\usepackage{color}
\usepackage[reqno]{amsmath}
\usepackage{amsthm}
\usepackage{fancyvrb}
\usepackage{amssymb,enumerate}
\usepackage[all]{xy}
\usepackage{endnotes}
\usepackage{lscape}
\newtheorem{com}{Comment}
\usepackage{float}
\usepackage{hyperref}
\newtheorem{lem} {Lemma}
\newtheorem{prop}{Proposition}
\newtheorem{thm}{Theorem}
\newtheorem{defn}{Definition}
\newtheorem{cor}{Corollary}
\newtheorem{obs}{Observation}
\usepackage[compact]{titlesec}
\usepackage{dcolumn}
\usepackage{tikz}
\usetikzlibrary{arrows}
\usepackage{multirow}
\usepackage{xcolor}
\newcolumntype{.}{D{.}{.}{-1}}
\newcolumntype{d}[1]{D{.}{.}{#1}}
\definecolor{light-gray}{gray}{0.65}
\usepackage{url}
\usepackage{listings}
\usepackage{color}

\definecolor{codegreen}{rgb}{0,0.6,0}
\definecolor{codegray}{rgb}{0.5,0.5,0.5}
\definecolor{codepurple}{rgb}{0.58,0,0.82}
\definecolor{backcolour}{rgb}{0.95,0.95,0.92}

\lstdefinestyle{mystyle}{
	backgroundcolor=\color{backcolour},   
	commentstyle=\color{codegreen},
	keywordstyle=\color{magenta},
	numberstyle=\tiny\color{codegray},
	stringstyle=\color{codepurple},
	basicstyle=\footnotesize,
	breakatwhitespace=false,         
	breaklines=true,                 
	captionpos=b,                    
	keepspaces=true,                 
	numbers=left,                    
	numbersep=5pt,                  
	showspaces=false,                
	showstringspaces=false,
	showtabs=false,                  
	tabsize=2
}
\lstset{style=mystyle}
\newcommand{\Sref}[1]{Section~\ref{#1}}
\newtheorem{hyp}{Hypothesis}

\title{Problem Set 1}
\date{Due: October 22, 2026}
\author{
	Applied Stats/Quant Methods 1\\
	Kanon Saito
}

\begin{document}
	\maketitle
	\section*{Instructions}
\begin{itemize}
	\item Please show your work! You may lose points by simply writing in the answer. If the problem requires you to execute commands in \texttt{R}, please include the code you used to get your answers. Please also include the \texttt{.R} file that contains your code. If you are not sure if work needs to be shown for a particular problem, please ask.
	\item Your homework should be submitted electronically on GitHub.
	\item This problem set is due before 23:59 on Thursday October 22, 2026. No late assignments will be accepted.

\end{itemize}

	
	\vspace{.5cm}
	\section*{Question 1: Political Science}
		\vspace{.25cm}
	The following table was created using the data from a study run in a major Latin American city.\footnote{Fried, Lagunes, and Venkataramani (2010). ``Corruption and Inequality at the Crossroad: A Multimethod Study of Bribery and Discrimination in Latin America. \textit{Latin American Research Review}. 45 (1): 76-97.} As part of the experimental treatment in the study, one employee of the research team was chosen to make illegal left turns across traffic to draw the attention of the police officers on shift. Two employee drivers were upper class, two were lower class drivers, and the identity of the driver was randomly assigned per encounter. The researchers were interested in whether officers were more or less likely to solicit a bribe from drivers depending on their class (officers use phrases like, ``We can solve this the easy way'' to draw a bribe). The table below shows the resulting data.

\newpage
\begin{table}[h!]
	\centering
	\begin{tabular}{l | c c c }
		& Not Stopped & Bribe requested & Stopped/given warning \\
		\\[-1.8ex] 
		\hline \\[-1.8ex]
		Upper class & 14 & 6 & 7 \\
		Lower class & 7 & 7 & 1 \\
		\hline
	\end{tabular}
\end{table}

\begin{enumerate}
	
	\item [(a)]
	Calculate the $\chi^2$ test statistic by hand/manually (even better if you can do "by hand" in \texttt{R}).\\
	R:
	\lstinputlisting[firstline=22, lastline=32]{PS02_answers_KS.R}
	\begin{verbatim}
	> chi_square
	[1] 3.791168
	\end{verbatim}

	\vspace{0.5cm}
	
	\item [(b)]
	Now calculate the p-value from the test statistic you just created (in \texttt{R}).\footnote{Remember frequency should be $>$ 5 for all cells, but let's calculate the p-value here anyway.}  What do you conclude if $\alpha = 0.1$?\\
	R:
	\lstinputlisting[firstline=36, lastline=36]{PS02_answers_KS.R}
	\begin{verbatim}
	> p_value
	[1] 0.1502306
	\end{verbatim}
		Since the p-value is greater than the threshold ($\alpha = 0.10$), we fail to reject the null hypothesis that driver class and bribe solicitation are independent.
	\newpage
	\item [(c)] Calculate the standardized residuals for each cell and put them in the table below.
	\vspace{1cm}
	
	\begin{table}[h]
		\centering
		\begin{tabular}{l | c c c }
			& Not Stopped & Bribe requested & Stopped/given warning \\
			\\[-1.8ex] 
			\hline \\[-1.8ex]
			Upper class  & 0.14 & -0.82 & 0.82 \\
			\\
			Lower class & -0.18 & 1.09  & -1.10  \\
		\end{tabular}
	\end{table}
	\lstinputlisting[firstline=39, lastline=45]{PS02_answers_KS.R}
	
	\vspace{1cm}
	\item [(d)] How might the standardized residuals help you interpret the results?  \\
	
	Standardized residuals help us interpret the results by showing the direction and magnitude of the differences between observed and expected values for each cell.
	\begin{itemize}
	\item Positive values indicate that a cell has more observations than expected under the null hypothesis.
	\item Negative values indicate that a cell has fewer observations than expected under the null hypothesis.
	\item Absolute values greater than 1.96 (or roughly 2.0) indicate a statistically significant deviation at the 5\% significance level.
	\end{itemize}
\end{enumerate}
\newpage

\section*{Question 2: Economics}
Chattopadhyay and Duflo were interested in whether women promote different policies than men.\footnote{Chattopadhyay and Duflo. (2004). ``Women as Policy Makers: Evidence from a Randomized Policy Experiment in India. \textit{Econometrica}. 72 (5), 1409-1443.} Answering this question with observational data is pretty difficult due to potential confounding problems (e.g. the districts that choose female politicians are likely to systematically differ in other aspects too). Hence, they exploit a randomized policy experiment in India, where since the mid-1990s, $\frac{1}{3}$ of village council heads have been randomly reserved for women. A subset of the data from West Bengal can be found at the following link: \url{https://raw.githubusercontent.com/kosukeimai/qss/master/PREDICTION/women.csv}\\

\noindent Each observation in the data set represents a village and there are two villages associated with one GP (i.e. a level of government is called "GP"). Figure~\ref{fig:women_desc} below shows the names and descriptions of the variables in the dataset. The authors hypothesize that female politicians are more likely to support policies female voters want. Researchers found that more women complain about the quality of drinking water than men. You need to estimate the effect of the reservation policy on the number of new or repaired drinking water facilities in the villages.
\vspace{.5cm}
\begin{figure}[h!]
	\caption{\footnotesize{Names and description of variables from Chattopadhyay and Duflo (2004).}}
	\vspace{.5cm}
	\centering
	\label{fig:women_desc}
	\includegraphics[width=0.9\textwidth]{/Users/kanon/AppliedStats1/StatsI_2026/problem_sets/PS02/my_answers/PS02_variables_table.png}
\end{figure}		

\newpage
\begin{enumerate}
	\item [(a)] State a null and alternative (two-tailed) hypothesis. \\
	
	$H_0$: The number of new or repaired drinking water facilities is the same between villages with female leaders and those with male leaders.\\
	
	$H_1$: The number of new or repaired drinking water facilities is different between villages with female leaders and those with male leaders.
	\vspace{2cm}
	\item [(b)] Run a bivariate regression to test this hypothesis in \texttt{R} (include your code!).\\
	R:
	\lstinputlisting[firstline=52, lastline=53]{PS02_answers_KS.R}
	\begin{verbatim}
	Call:
	lm(formula = water ~ female, data = women_data)
	
	Residuals:
	Min     1Q Median     3Q    Max 
	-22.68 -14.78  -7.81   2.29 317.32 
	
	Coefficients:
	Estimate Std. Error t value Pr(>|t|)    
	(Intercept)   14.813      2.382   6.220 1.56e-09 ***
	female         7.864      3.838   2.049   0.0413 *  
	---
	Signif. codes:  0 ‘***’ 0.001 ‘**’ 0.01 ‘*’ 0.05 ‘.’ 0.1 ‘ ’ 1
	
	Residual standard error: 33.51 on 320 degrees of freedom
	Multiple R-squared:  0.01295,	Adjusted R-squared:  0.009867 
	F-statistic: 4.199 on 1 and 320 DF,  p-value: 0.04126
	\end{verbatim}
	\vspace{6cm}
	\item [(c)] Interpret the coefficient estimate for reservation policy. \\
	\begin{itemize}
	\item The intercept ($\alpha$) indicates that when $X = 0$, the expected value of $Y$ is 14.813.
	\item The slope ($\beta$) indicates that for every one-unit increase in $X$, the value of $Y$ increases by 7.864.
	\item The $p$-value is 0.0413 ($p < 0.05$), which means we can reject the null hypothesis at the 95\% confidence level that the number of new or repaired drinking water facilities is the same between villages with female leaders and those with male leaders.
	\end{itemize}
\end{enumerate}


%\end{enumerate}
\end{document}
